\subsection*{The full dihedral centralizer and its parameter orbifold}

\statement{B.74}{Proved}{A canonical direct factor and its full profinite comparison}
Let \(D=\langle d\rangle\simeq C_5\) and
\(J=\langle d,r\rangle\simeq D_5\) be B.65's actual groups,
and write \(\Gamma=\Gamma_{0,[12]}\).
Then
\[
 \begin{split}
 C_\Gamma(D)&=D\times C_\Gamma(J),\\
 C_\Gamma(J)&\simeq C_2*\mathbb Z,\\
 C_{\widehat\Gamma}(J)
       &=\widehat{C_\Gamma(J)}
          \simeq\widehat{C_2*\mathbb Z}.                         
 \end{split}                                                    \tag{B.74.1}
\]
The first splitting is canonical for the specified \(r\).
Its pure subgroup is \(H_5\), normal of index ten in
\(C_\Gamma(J)\), with quotient \(D_5\).
The associated actual orbifold covering is
\[
 \mathbb P^1_u\setminus\{0,\infty,\mu_5\}
       \longrightarrow
 \bigl(\mathbb P^1_q\setminus\{4,\infty\};\
                     \text{cone order }2\text{ at }0\bigr),
 \qquad q=\frac{(u^5+1)^2}{u^5}.                               \tag{B.74.2}
\]

\begin{proof}
Put \(C=C_\Gamma(D)\).  The two fixed punctures \(p,q_0\)
of \(D\) are individually preserved by \(C\), since their
local rotation residues are \(+1,-1\pmod5\).
The two free \(D\)-orbits \(A,B\) may be interchanged.
The cyclic-cover comparison gives
\[
        1\longrightarrow D\longrightarrow C
           \longrightarrow Q_4\longrightarrow1,               \tag{B.74.3}
\]
where \(Q_4\) fixes \(\bar p,\bar q_0\) individually and
preserves the unordered pair \(\{\bar A,\bar B\}\).
Every such base class lifts: it preserves the cyclic
covering character with peripheral values
\(d,d^{-1},1,1\), respectively, and induces the identity
on \(D\).  This proves the indicated surjectivity.

The image \(\bar r\) in the full four-point mapping-class
group is the Klein-four double transposition
\[
                 (\bar p\ \bar q_0)(\bar A\ \bar B).
\]
The pure four-point group centralizes the canonical
Klein four group.  The only additional puncture
permutation allowed in \(Q_4\) is
\((\bar A\ \bar B)\), which also fixes this
double transposition under conjugation.
Thus \(Q_4\) centralizes \(\bar r\), although \(\bar r\)
itself interchanges the two pole labels and is not an
element of \(Q_4\).
It follows that
\[
                 \delta(c)=[c,r]\in D\qquad(c\in C).             \tag{B.74.4}
\]
We use \([c,r]=crc^{-1}r^{-1}\).
Because \(D\) is central in \(C\), the function
\(\delta:C\to D\) is a homomorphism.
Since \(rdr^{-1}=d^{-1}\), its restriction to \(D\)
is the squaring automorphism.
There is consequently a homomorphic retraction
\[
             \rho:C\longrightarrow C_\Gamma(J),
             \qquad
             \rho(c)=\delta(c)^{-1/2}c.                         \tag{B.74.5}
\]
The half power is unique in \(C_5\), explicitly the
third power.  The kernel of \(\rho\) is \(D\),
and \(D\cap C_\Gamma(J)=1\).
This proves the actual direct product in (B.74.1);
it also identifies \(C_\Gamma(J)\) with \(Q_4\).
This complement is distinguished by \(r\).
It is not silently identified with B.52's separately
normalized vertex-fixing rank-two complement.

Normalize the ordered four-point configuration to
\(\{0,\infty,1,t\}\).
Interchanging the two free orbit labels sends \(t\)
to \(t^{-1}\).  The quotient coordinate is
\[
                         q=t+2+t^{-1}.                         \tag{B.74.6}
\]
The fixed point \(t=-1\) gives a cone point of order
two at \(q=0\).
The deleted point \(t=1\) gives the cusp \(q=4\);
the two deleted points \(t=0,\infty\) give the cusp
\(q=\infty\).
Thus \(Q_4\) is the fundamental group of an orbifold
with two cusps and one order-two cone point, namely
\(C_2*\mathbb Z\).
In the actual ordered cyclic family \(t=u^5\),
which gives (B.74.2).
Its degree and local types are
\[
\begin{array}{c|c|c}
\text{target}&\text{source points}&\text{local index}\\\hline
0&u^5=-1\quad(5\text{ points})&2\\
4&u^5=1\quad(5\text{ deleted points})&2\\
\infty&0,\infty\quad(2\text{ deleted points})&5 .
\end{array}                                                     \tag{B.74.7}
\]
The deck transformations \(u\mapsto\zeta u\) and
\(u\mapsto u^{-1}\) generate \(D_5\).
The cover is unramified as a cover of the displayed
orbifolds.  Therefore its source group \(H_5\)
is normal of index ten, with the stated quotient.
For orbifold meridians \(e\) at zero and \(a\) at four,
the quotient can be written
\[
 \langle e,a\mid e^2=1\rangle
       \longrightarrow
 \langle e,a\mid e^2=a^2=(ea)^5=1\rangle\simeq D_5.              \tag{B.74.8}
\]
All these maps use the based source curve and common
marking already specified in B.28 and B.61.

For the full profinite comparison, the valid
prime-order-five theorem used in B.28 identifies
\(C_{\widehat\Gamma}(D)\) with \(\widehat C\).
Complete the proved direct product:
\[
          C_{\widehat\Gamma}(D)
                    =D\times\widehat{C_\Gamma(J)}.
\]
Conjugation by \(r\) inverts \(D\) and is the identity
on the second factor, by its very definition.
Since the fixed subgroup of inversion in \(C_5\)
is trivial, taking \(r\)-fixed points gives the
last assertion of (B.74.1).
This argument does not exchange uncontrolled
completed intersections and requires no
nonsolvable centralizer assertion.
\end{proof}

For B.73's smaller ambient realization, let
\(\widetilde Q_8\) fix the ordinary label from
\(\{1,3\}\), permit all permutations of the other
five ordinary labels, and preserve the branch pair.
It contains \(Q_8\) with index \(120\).
Choosing the lift that fixes original label \(1\)
again splits its deck extension.
The image of \(C_\Gamma(J)\) in
\(\widetilde Q_8\) is faithful, including after
completion: the chosen section is an open subgroup
of the actual involution centralizer, whose full
topology is the valid prime-order-two comparison.
Equivalently the only possible deck kernel
\(\langle r\rangle\) has trivial intersection with
the \(J\)-centralizer.

On the ten non-pole original labels, \(J\) acts regularly.
Its permutation centralizer is the right regular \(D_5\),
and the index-ten calculation shows that
\(C_\Gamma(J)\) realizes all of it.
Its action on the five \(r\)-orbits is the faithful
degree-five action on \(\langle r\rangle\backslash J\);
the core of \(\langle r\rangle\) in \(D_5\) is trivial.
Consequently
\[
       C_\Gamma(J)\cap Q_8=H_5
\]
under this faithful identification.
The full \(C_2*\mathbb Z\) therefore belongs in
\(\widetilde Q_8\), not in its ordinary-label-fixing
subgroup \(Q_8\).
Finally, (B.74.5) is a retraction with domain
\(C_\Gamma(D)\), which has infinite index in the
ambient mapping-class group.  It supplies no
finite-index ambient retraction and leaves the
refined-product problem open.

